% titel: Blatt 6, A1: Komplexe Zahlen
% skript: Kap. 1 (Komplexe Zahlen, Polarkoordinaten)

% Teil I: Werkzeuge, kein Lösungsweg.

Polarkoordinaten einer komplexen Zahl bestimmen, eine Kreisscheibe in der
komplexen Ebene charakterisieren und den dritten Eckpunkt eines
gleichseitigen Dreiecks finden.

\begin{werkzeugkasten}[Polarkoordinaten]
Eine komplexe Zahl $z=a+ib$ schreibt sich mit Radius $r$ und Winkel $\varphi$
als $z=r(\cos\varphi+i\sin\varphi)$ mit
\begin{equation*}
  r=|z|=\sqrt{a^2+b^2}, \qquad
  \tan\varphi=\frac{b}{a}\ \text{(im richtigen Quadranten)}.
\end{equation*}
Arithmetik wird im Polarbild einfach:
\begin{equation*}
  z_1\cdot z_2=r_1r_2\,\mathrm e^{i(\varphi_1+\varphi_2)},\qquad
  z^n=r^n\,\mathrm e^{in\varphi}.
\end{equation*}
Für $z=1+i\sqrt3$: $r=\sqrt{1+3}=2$, $\varphi=\arctan(\sqrt3)=\frac\pi3$.

\textit{Skript:} Kap. 1 (Polarkoordinaten).
\end{werkzeugkasten}

\begin{werkzeugkasten}[Mengen und Abstände in $\mathbb C$]
$|z-z_0|$ ist der Abstand von $z$ zum Punkt $z_0$. Daher beschreibt
\begin{itemize}
  \item $|z-z_0|=r$ den Kreis um $z_0$ mit Radius $r$,
  \item $|z-z_0|<r$ den offenen Kreis~\emph{innen} (ohne Rand),
  \item $|z-z_0|>r$ das Gebiet~\emph{au{\ss}erhalb} des Kreises.
\end{itemize}

\textit{Skript:} Kap. 1 (Komplexe Zahlen).
\end{werkzeugkasten}

\begin{achtung}[Darauf achten]
\begin{itemize}
  \item $\arctan(b/a)$ liefert den Winkel nur in $(-90^\circ,90^\circ)$;
        den Quadranten von $z$ separat beachten (hier liegt $1+i\sqrt3$ im
        ersten Quadranten, $\varphi=\pi/3$ bleibt).
  \item $|z-2i|<2$ ist ein \emph{offener} Kreis: der Rand
        ($|z-2i|=2$) gehört nicht dazu.
  \item Bei zwei gegebenen Eckpunkten eines gleichseitigen Dreiecks gibt es
        zwei mögliche dritte Punkte (die beiden Seiten-Einheits-Richtungen).
\end{itemize}
\end{achtung}
